\section{Target Debit and Exact-Fallback Control}
\label{sec:theory}

We first transfer regret from delivered rows to the dense target game, then
turn the resulting correction into a causal plan constraint.

Define delivered weighted positive local regret
\begin{equation}
\widehat{\mathcal R}_T^w=
\sum_\rho\pos{\max_a\sum_{t=1}^T w_t h_{\rho,t}(a)}.
\label{eq:delivered-positive-regret}
\end{equation}
Let
$\Omega_0=\max_{x,x',y}\{u(x,y)-u(x',y)\}$, which is finite on the
realization-plan polytopes.

\begin{theorem}[Weighted target-debit transfer]
\label{thm:weighted-target-debit}
For every finite event-faithful execution under the declared additive public
cut contract,
\begin{equation}
\boxed{
\operatorname{expl}(\bar x_T^w,\bar y_T^w)
\le\frac{1}{2W_T}\left[
\widehat{\mathcal R}_T^w
+\sum_\rho\psi\!\left(\sum_{t=1}^T w_t\epsilon_{\rho,t}\right)
+\Omega_0w_T+\sum_{t=1}^T w_td_t
\right].}
\label{eq:weighted-target-debit}
\end{equation}
\end{theorem}

\paragraph{Proof.}
Put $H_\rho=\sum_t w_th_{\rho,t}$ and
$E_\rho=\sum_t w_t\epsilon_{\rho,t}$, so the cumulative target row is
$H_\rho-E_\rho$.  Rowwise,
\begin{equation}
\pos{\max_a(H_\rho(a)-E_\rho(a))}
\le \pos{\max_a H_\rho(a)}+\psi(E_\rho).
\label{eq:one-sided-row-transfer}
\end{equation}
Indeed, for every action $a$,
$H_\rho(a)-E_\rho(a)\le\max_bH_\rho(b)+\max_b[-E_\rho(b)]$;
taking the maximum and then the positive part proves the display.  Summing it
over the finite row registry bounds the two players' exact target local
regrets by $\widehat{\mathcal R}_T^w+\sum_\rho\psi(E_\rho)$.

For completeness, fix a valid comparator for one player.  The perfect-recall
policy-difference identity writes its weighted external regret as a sum over
owned information sets of the comparator's own reach, which lies in $[0,1]$,
times a comparator mixture of the cumulative centered counterfactual row.
That mixture is at most the row maximum, and replacing that maximum by its
positive part preserves the bound.  Applying this argument at
$(x_t,y_t)$ for player zero and at $(x_{t+1},y_t)$ for player one, then taking
the two comparator suprema, gives the native alternating local-to-external
reduction.  Bilinearity of $u$ identifies its weighted comparator sums with
the gap of the normalized weighted realization-plan averages.  The only
unmatched term is
\begin{align}
C_T^w={}&\max_y\sum_t w_tu(x_{t+1}-x_t,y)
-\sum_t w_tu(x_{t+1}-x_t,y_t). \label{eq:weighted-bridge}
\end{align}
Summation by parts gives
\[
\sum_t w_t(x_{t+1}-x_t)=w_Tx_{T+1}-w_1x_1
-\sum_{t=2}^T(w_t-w_{t-1})x_t.
\]
The negative endpoint coefficient has magnitude $w_1$, while the increment
coefficients have total $w_T-w_1$; together their magnitude is $w_T$.
Because $W_T>0$ and the clock is nonnegative and nondecreasing, $w_T>0$.
The negative terms divided by $w_T$ therefore form a convex combination
$x^-$ of feasible realization plans, and, for every $y$,
\[
\sum_t w_tu(x_{t+1}-x_t,y)
=w_T\{u(x_{T+1},y)-u(x^-,y)\}\le \Omega_0w_T.
\]
For the second term,
$-u(x_{t+1}-x_t,y_t)=u(x_t,y_t)-u(x_{t+1},y_t)\le d_t$.
Thus $C_T^w\le\Omega_0w_T+\sum_tw_td_t$.  Substituting the rowwise transfer
into the native reduction and dividing the Nash gap by two proves the
theorem. \hfill$\square$

\paragraph{Why the correction is one-sided.}
The learner accumulates $H_\rho$, whereas the target accumulates
$H_\rho-E_\rho$.  For a two-action row with $H_\rho=(r,0)$, positive error
$E_\rho=(e,0)$ with $e\ge0$ cannot increase the target maximum; negative
error $E_\rho=(-e,0)$ can increase it by $e$.  Thus only cumulative
underestimation can create additional positive target regret, exactly as
measured by $\psi(E_\rho)$.  The four numerator terms in
Equation~\ref{eq:weighted-target-debit} have separate roles: delivered learner
regret, public target debit, the endpoint cost of the weighted alternating
observer, and payoff decrease during player-zero updates.  A symmetric norm
is valid but discards the first distinction.  The appendix gives
component-resolved identities that compute or upper-bound the same debit from
retained traces.

\begin{proposition}[Complete coverage is a zero-debit action]
\label{prop:zero-debit-action}
For additive signed RM or RM+ with the uniform zero-positive-mass fallback,
one event of complete, same-origin public
coverage under the event-faithful simultaneous-update contract satisfies
$\epsilon_{\rho,t}=0$ for every row and $d_t=0$.
\end{proposition}

Complete coverage has $B=N$, so its multiplier is one and
$h_{\rho,t}=g_{\rho,t}$ at the same immutable origin; hence
$\epsilon_{\rho,t}=0$.  For the payoff term, consider one updated row and put
$p=\operatorname{RM}(z)$, $h=v-\langle p,v\rangle\one$,
$x=[z]_+$, and $y=[z+h]_+$.  Coordinatewise monotonicity gives
$(y_a-x_a)h_a\ge0$.  If $x$ has positive mass, it is a positive multiple of
$p$; otherwise $x=0$.  In either case $x^Th=0$, hence $y^Th\ge0$.
If $y$ has positive mass, normalization gives
$\langle\operatorname{RM}(z+h),h\rangle\ge0$.  If $y$ has zero mass, the old
mass must also be zero (a positive-mass centered state cannot be projected to
zero), so both policies use the same uniform fallback and the inner product
is zero.  Since $h$ is centered at $p$, this proves
$\langle p^+-p,v\rangle\ge0$ in every mass branch.  RM+ has the same next
policy because applying positive part once more is policy-idempotent.

Let $\sigma$ be the immutable old profile and $\sigma^+$ the simultaneous
update of all rows owned by player zero; the opponent rows agree.  Perfect
recall gives the exact identity
\[
u(\sigma^+,y_t)-u(\sigma,y_t)
=\sum_{I\in\mathcal I_0}\pi_0^{\sigma^+}(I)
 \sum_a(\sigma^+(I,a)-\sigma(I,a))v^\sigma(I,a).
\]
It follows by telescoping terminal-path probabilities and grouping each
changed decision occurrence by its information set; perfect recall makes the
new-policy own reach $\pi_0^{\sigma^+}(I)$ independent of the representative
history.  Every reach factor is nonnegative and every row inner product is
nonnegative by the preceding argument.  Therefore
$u(x_{t+1},y_t)\ge u(x_t,y_t)$, so $d_t=0$.

\subsection{A debit-budgeted controller}

Let $T_k=km$ be aligned checkpoints.  For the executed plan in block $k$,
define its realized weighted charge
\begin{equation}
q_k=
\sum_\rho\psi\!\left(\sum_{t=T_{k-1}+1}^{T_k}
w_t\epsilon_{\rho,t}\right)
+\sum_{t=T_{k-1}+1}^{T_k}w_td_t.
\label{eq:block-charge}
\end{equation}
A partial block has a history-measurable finite envelope $\bar q_k$ satisfying
$0\le q_k\le\bar q_k$; its realized charge may be revealed after execution.
The envelope enters the guarantee and need not be evaluated by the planner.
Complete coverage has $q_k=0$ by
Proposition~\ref{prop:zero-debit-action} and settles immediately.

Let $A_{k-1}$ be the sum of settled partial charges and let $P_{k-1}$ contain
the completed partial blocks whose charges remain unsettled.  Thus the
mathematical realized account is
$Q_{k-1}=A_{k-1}+\sum_{j\in P_{k-1}}q_j$ even when the controller has not yet
observed every term.  Fix a predetermined nonnegative nondecreasing budget
$b(T_k)$ and an outstanding-audit cap $\ell\ge1$.  DB-CFR may select a partial
plan only if
\begin{equation}
A_{k-1}\le b(T_k),\qquad |P_{k-1}|<\ell;
\label{eq:audited-admission}
\end{equation}
otherwise it executes complete coverage.  Within the solvent branch, any
fixed completed-history rule may choose width and order.  Let
$R_k^{(\ell)}$ be the sum of the $\ell$ largest envelopes among partial blocks
executed through $k$, padding with zeros.

\begin{theorem}[Causal exact-fallback guarantee]
\label{thm:debit-controller}
Assume event-faithful candidate executions; a designated complete-coverage
candidate executable from every reachable block boundary; a finite
history-measurable envelope for every selectable partial candidate; globally
zero-initialized additive RM+; centered delivered rows bounded uniformly over
the finite candidate registry; and a fixed nondecreasing output clock.  For
every legal audit settlement order, DB-CFR is well defined,
$Q_K\le b(T_K)+R_K^{(\ell)}$, and
\begin{equation}
\boxed{
\operatorname{expl}(\bar x_{T_K}^w,\bar y_{T_K}^w)
\le
\frac{C_Rw_{T_K}\sqrt{T_K}+\Omega_0w_{T_K}+b(T_K)
+R_K^{(\ell)}}
{2W_{T_K}}.}
\label{eq:controller-bound}
\end{equation}
Under uniform weights, the same statement holds for additive signed RM.
\end{theorem}

\paragraph{Proof.}
The executable complete candidate makes the transition total whenever the
partial admission test fails; otherwise the selected partial candidate and
its finite envelope determine the next event-faithful block.  Thus the stated
controller is well defined under the theorem's premises.

Suppose block $k$ is partial.  Equation~\ref{eq:audited-admission} gives
$A_{k-1}\le b(T_k)$.  Its prior pending blocks together with the new block
form a subset of the executed partial blocks of cardinality at most $\ell$.
Since $q_j\le\bar q_j$ for each member,
\[
Q_k=A_{k-1}+\sum_{j\in P_k}q_j
\le b(T_k)+\sum_{j\in P_k}\bar q_j
\le b(T_k)+R_k^{(\ell)}.
\]
If block $k$ is complete, Proposition~\ref{prop:zero-debit-action} gives
$q_k=0$ and the realized account is unchanged.  Monotonicity of $b$ and of
the top-$\ell$ reserve closes this branch of the induction.  Moving a charge
from pending to settled adds it to $A_k$ and removes the same amount from the
pending sum, so arbitrary audit settlement leaves $Q_k$ invariant.

For each row, $\psi(z)=\max_a[-z(a)]_+$ is subadditive.  Therefore
\[
\sum_\rho\psi\!\left(\sum_kE_{\rho,k}\right)
\le\sum_k\sum_\rho\psi(E_{\rho,k}),
\]
and the nonnegative payoff-decrease charges add exactly.  Hence the two debit
terms in Theorem~\ref{thm:weighted-target-debit} are at most $Q_K$.
Finally, the coordinatewise RM+ clipping telescope gives, for each row with
$\lVert h_{\rho,t}\rVert_2\le G_\rho$,
$[\max_a\sum_t w_th_{\rho,t}(a)]_+
\le w_{T_K}G_\rho\sqrt{T_K}$.  Summing the finite row registry and taking
$C_R=\sum_\rho G_\rho$ yields the stated learner term.  Substitute
$Q_K\le b(T_K)+R_K^{(\ell)}$ into
Theorem~\ref{thm:weighted-target-debit}.  For uniform weights, the separate
positive-part potential argument for unclipped signed RM supplies the same
$C_R\sqrt{T_K}$ form. \hfill$\square$

For a fixed block length, bounded row errors give a finite constant $C_B$ with
$\bar q_k\le C_B\sum_{t=T_{k-1}+1}^{T_k}w_t$.  If $w_t=t^p$, $p\ge0$, and
$b(T)=\kappa w_T\sqrt T$, then
$R_K^{(\ell)}\le\ell C_Bmw_{T_K}$ and
\[
\operatorname{expl}\le\frac{p+1}{2}\left[
\frac{C_R+\kappa}{\sqrt T}
+\frac{\Omega_0+\ell C_Bm}{T}\right].
\]
Thus fixed $\ell$ contributes only $O(T^{-1})$ and preserves the additive-RM+
rate.  The one-block post-execution controller is the case $\ell=1$.  If a
continuation-valid charge is instead available before selection, requiring
$Q_{k-1}+q_k\le b(T_k)$ recovers the zero-overdraft prospective controller as
a special implementation.

\paragraph{Comparison with symmetric charging.}
For every completed block, one-sided debit is no larger than charging
$\sum_\rho\lVert E_\rho\rVert_\infty$ plus the same payoff term.  The
inclusion is strict: the safe plan in the witness below has one-sided charge
zero and symmetric charge $7/4$.  The appendix gives the formal dominance
statement and proof.

\subsection{Why counts cannot certify the decision}

\begin{proposition}[Equal-metadata debit separation]
\label{prop:metadata-separation}
There is a finite two-player zero-sum perfect-recall game and two four-round
singleton public-cut plans $S^+$ and $S^-$ with identical public counts, two
complete coverage cycles, batch width, symbolic work and task graph, maximum
prefix discrepancy $D_4=1$, and scalar strategy-path length $P_4=4$, but
\begin{equation}
q(S^+)=0,\qquad q(S^-)=\frac{11}{8}.
\label{eq:charge-separation}
\end{equation}
Consequently, any charge based only on these shared summaries and sound for
both plans assigns both at least $11/8$; at budget one it rejects $S^+$ even
though its exact debit is zero.  Likewise, a pathwise strict-budget admission
rule restricted to these summaries must reject both before execution.  To
admit the zero-debit plan, a controller must obtain semantic information
before execution or permit a bounded post-execution correction.
\end{proposition}

The witness starts zero-state signed RM on the two public conditional matrices
\[
M_1=\begin{pmatrix}3&3\\1&1\end{pmatrix},\qquad
M_2=\begin{pmatrix}-1&0\\2&-1\end{pmatrix},
\]
and compares the balanced orders $S^+=(1,2,1,2)$ and
$S^-=(2,1,2,1)$.  Exact replay gives
\begin{center}
\small
\begin{tabular}{@{}lcccc@{}}
\toprule
Plan & player-zero error & player-one error & payoff charge & $q$ \\
\midrule
$S^+$ & $(3/4,3/4)$ & $(0,1)$ & $0$ & $0$ \\
$S^-$ & $(-3/8,5/8)$ & $(-3/4,3/4)$ & $1/4$ & $11/8$ \\
\bottomrule
\end{tabular}
\end{center}
Both error rows of $S^+$ are coordinatewise nonnegative.  Plan $S^-$ instead
incurs row underestimation $3/8+3/4$ and payoff decrease $1/4$.  The appendix
gives the full recurrence and rational verification; its symmetric total
charges are $7/4$ and $13/8$, respectively.  The result establishes a finite
decision separation for the declared summary class.
